\begin{abstract}
Planning uncrewed aerial vehicle (UAV) relief flights after a flood requires an operator to read a flood map, choose delivery points, route around water and author an autopilot mission by hand. Prior work seldom produces a mission that an autopilot can load. This paper presents \adapt{}, an open pipeline that converts a flood-annotated map raster and current weather into QGroundControl WPL~110 missions. It combines operator-sampled HSV segmentation, a weather-parameterised flood-spread rule, per-region drop points, nearest-neighbour ordering, per-leg D*~Lite planning and a checked pixel-to-geodetic conversion. It supports two approaches: one drone visits every drop point, or a fleet of up to $K$ drones is used, in which each drone serves the cluster of drop points nearest its own base under a battery model. Fleet plans whose routes cross are rejected. On three published flood maps with a deterministic simulated operator, every generated mission passes an independent validator. Randomised coordinate round trips agree within $7\times10^{-8}$~px, and 31 of 33 injected defects are detected. Raising $K$ from 1 to 4 increases the drops served from 10 to 18 of 22 and from 17 to 34 of 41 under moderate weather. Limitations remain. Predicted-flood obstacles cause all 14 no-path legs of the single-UAV mode, one thematic map is segmented as 83\% flood, the multi-base assignment does not balance workload, and the maps are not georeferenced. No ground-station, simulation or flight validation has been performed, so \adapt{} is a verified research baseline, not a deployable system.
\end{abstract}

\begin{IEEEkeywords}
Unmanned aerial vehicles, disaster response, flood mapping, mission planning, multi-UAV planning, path planning, D* Lite, geospatial transformation, software verification.
\end{IEEEkeywords}
